\chapter*{Résumé}
\addcontentsline{toc}{chapter}{Résumé}

Ce travail a été effectué dans le cadre de la réalisation de notre projet de fin d'études au sein de \textbf{Loomens}, une agence tunisienne de développement web et de transformation digitale située au Cyberparc de Djerba, en vue de l'obtention du diplôme de Master Professionnel en Génie Logiciel, décerné par l'Institut Supérieur d'Informatique et de Mathématiques de Monastir (ISIMM).

\medskip

La plateforme \textbf{LoomStay} permet aux établissements hôteliers de digitaliser plusieurs services destinés aux clients et au personnel. Elle couvre notamment le check-in numérique avec fiches voyageurs, l'accès client à un espace chambre via QR code installable en tant que PWA, la gestion des réclamations avec classification automatique par intelligence artificielle, la communication multilingue avec traduction automatique, l'affectation des interventions aux employés de maintenance et de housekeeping, le suivi des tâches quotidiennes de ménage et la validation des interventions par scan QR à travers une application mobile Flutter.

\medskip

La solution repose sur une architecture orientée services composée d'une application web React/TypeScript, d'un backend Node.js/Express avec Prisma ORM et PostgreSQL, d'un service d'intelligence artificielle FastAPI intégrant le modèle pré-entraîné NLLB-200 pour la traduction et un classificateur TF-IDF / Régression Logistique pour la catégorisation des réclamations, ainsi que d'une application mobile Flutter pour les employés terrain.

\medskip

\textbf{Mots clés :} Hôtellerie, Digitalisation, PWA, React, Node.js, Flutter, Intelligence Artificielle, Machine Learning, Traduction automatique, NLLB-200, QR Code, Scrum.
\newpage
